% ============================================================================
%  Machine Learning Capstone Project - Summary Report
%  Self contained: no external images, no uncommon packages.
%  Compile with: pdflatex Project_Summary_Report.tex   (run twice for the ToC)
% ============================================================================
\documentclass[11pt,a4paper]{article}

\usepackage[T1]{fontenc}
\usepackage[utf8]{inputenc}
\usepackage[margin=2.4cm]{geometry}
\usepackage{booktabs}
\usepackage{array}
\usepackage{xcolor}
\usepackage{titlesec}
\usepackage{enumitem}
\usepackage[hidelinks]{hyperref}

\definecolor{deep}{HTML}{522CBA}
\definecolor{lilac}{HTML}{9B78F7}

\titleformat{\section}{\large\bfseries\color{deep}}{\thesection}{0.7em}{}
\titleformat{\subsection}{\normalsize\bfseries\color{deep}}{\thesubsection}{0.7em}{}
\titlespacing*{\section}{0pt}{1.4ex plus .3ex}{0.7ex}
\titlespacing*{\subsection}{0pt}{1.1ex plus .2ex}{0.5ex}

\setlist[itemize]{leftmargin=1.3em,topsep=2pt,itemsep=1.5pt,parsep=0pt}
\renewcommand{\arraystretch}{1.15}

\newcommand{\metric}[1]{\textbf{\color{deep}#1}}

\begin{document}

% ---------------------------------------------------------------- title block
\begin{center}
    {\large\color{lilac}\bfseries MACHINE LEARNING CAPSTONE PROJECT}\\[0.8em]
    {\LARGE\bfseries\color{deep} Predicting Flight Arrival Delay}\\[0.35em]
    {\LARGE\bfseries\color{deep} from Post Departure Information}\\[1.0em]
    {\normalsize Khalid Al Dosari \quad\textbullet\quad Abdulaziz Alshareef
     \quad\textbullet\quad Anas Alzahrani \quad\textbullet\quad Feras Madkhali}\\[0.4em]
    {\small Tuwaiq Bootcamp}
\end{center}
\vspace{0.6em}
\hrule
\vspace{1.2em}

% ---------------------------------------------------------------- 1
\section{Problem and Dataset}

\textbf{Objective.} Predict how late a US domestic flight will arrive, at the moment it
pushes back from the gate. Two targets are modelled on the same rows, the same features
and the same split: \texttt{ArrDel15} (classification, does the flight land 15 minutes
late or more) and \texttt{ArrDelay} (regression, arrival delay in minutes, which may be
negative when a flight arrives early).

\textbf{Real world motivation.} Once an aircraft has left the gate, airlines, airports
and connecting passengers need an arrival estimate sharper than the printed schedule.
Gate assignment, crew rotation and passenger connection decisions all depend on it.

\textbf{Dataset.} Kaggle, \emph{Flight Delay Dataset 2018--2022}, file
\texttt{Combined\_Flights\_2018.csv}, originally published by the US Department of
Transportation, Bureau of Transportation Statistics (On Time Performance database).

\begin{center}
\small\url{https://www.kaggle.com/datasets/robikscube/flight-delay-dataset-20182022}
\end{center}

One row is one scheduled domestic flight. The raw file contains \metric{5{,}689{,}512
rows} and \metric{61 columns}, mixing numerical measurements, categorical identifiers
and schedule fields.

\textbf{Why machine learning is appropriate.} Arrival delay depends on how late the
aircraft already is, ground congestion, route, carrier and time of day, all interacting.
No closed form rule captures this, but a model can learn the pattern from millions of
historical flights.

% ---------------------------------------------------------------- 2
\section{Data Quality and Cleaning}

Seven checks were run on the raw data. The decision taken mattered more than the
technique used, so each is justified rather than applied mechanically.

\begin{center}
\small
\begin{tabular}{@{}p{3.5cm} p{5.3cm} p{5.8cm}@{}}
\toprule
\textbf{Check} & \textbf{What was found} & \textbf{Decision and reason} \\
\midrule
Missing values & 1.8\% of arrival fields, 1.5\% of departure fields &
Drop cancelled (1.55\%) and diverted (0.25\%) flights. A flight that never arrived has
no arrival delay to predict, so imputing would invent data. \\
Duplicate rows & None found & Nothing to remove. \\
Incorrect types & \texttt{CRSDepTime} is an integer but encodes a clock time (1259 to
1300 is one minute, not 41) &
Convert to minutes since midnight, then to cyclic sine and cosine so 23:59 and 00:01 sit
close together. \\
Invalid values & 67 rows with negative scheduled or actual duration &
Drop. Physically impossible and unrecoverable; 67 of 5.69M is negligible. \\
Inconsistent categories & \texttt{Airline}, \texttt{Origin}, \texttt{Dest} checked for
case and spelling variants, none collide &
No cleaning required. \\
Outliers & IQR flags 9.4\% of \texttt{ArrDelay} and 13.1\% of \texttt{DepDelay};
skewness 8.5 and 9.6 &
\textbf{Keep all.} The long right tail is the phenomenon the project exists to predict.
Removing it would bias the model toward underestimating real delays. \\
Class imbalance & 80.5\% on time against 19.5\% delayed &
Report ROC-AUC, Precision, Recall and F1 rather than Accuracy alone; use
\texttt{class\_weight="balanced"} where supported. \\
\bottomrule
\end{tabular}
\end{center}

After cleaning, \metric{5{,}578{,}560 rows} remain.

% ---------------------------------------------------------------- 3
\section{Exploratory Analysis: Five Findings}

\begin{enumerate}[leftmargin=1.5em,itemsep=3pt]
  \item \textbf{Departure delay dominates.} \texttt{DepDelay} correlates with
        \texttt{ArrDelay} at \metric{0.96}. Once a flight leaves late it lands late
        almost as a rule. This single relationship sets both the ceiling and the floor
        for the entire project.
  \item \textbf{Delay compounds through the day.} Mean arrival delay rises from about
        $-3$ minutes for 5am departures to $+12$ minutes around 6pm, as aircraft
        rotations fall progressively behind schedule. Scheduled hour is therefore a
        cheap and strong feature.
  \item \textbf{Delay is seasonal.} Summer (June to August, 9.7 to 11.5 minutes mean)
        is roughly three times worse than autumn (September to October, about 3
        minutes), consistent with thunderstorm season. Calendar month carries real
        signal.
  \item \textbf{Carriers differ substantially at the same scale.} Among the ten largest
        carriers, JetBlue is delayed 27.3\% of the time against Delta at 12.5\%.
        Operating practice and hub congestion matter, justifying \texttt{Airline} as a
        feature despite being categorical.
  \item \textbf{Ground congestion adds independent signal.} \texttt{TaxiOut} correlates
        with \texttt{ArrDelay} at 0.22, moderate but real, and is not visible in
        \texttt{DepDelay} alone. Conversely \texttt{Distance} and
        \texttt{CRSElapsedTime} correlate at 0.98 with each other but only $-0.02$ with
        the target, so trip length alone says little about lateness and only one of the
        two is retained.
\end{enumerate}

% ---------------------------------------------------------------- 4
\section{Feature Preparation and the Leakage Boundary}

The single most consequential decision in this project is \emph{when} the prediction is
made. We fix that moment at departure, which determines exactly which columns are legal.

\begin{center}
\small
\begin{tabular}{@{}p{7.0cm} p{7.6cm}@{}}
\toprule
\textbf{Legal at departure (used)} & \textbf{Banned, known only after landing} \\
\midrule
\texttt{DepDelay}, \texttt{TaxiOut}, \texttt{CRSElapsedTime},
\texttt{dep\_sin}, \texttt{dep\_cos}, \texttt{Month}, \texttt{DayofMonth},
\texttt{DayOfWeek}, \texttt{is\_weekend}, \texttt{Airline}, \texttt{Origin},
\texttt{Dest}
&
\texttt{ArrTime}, \texttt{WheelsOn}, \texttt{TaxiIn},
\texttt{ActualElapsedTime}, \texttt{AirTime}, \texttt{ArrDelayMinutes},
\texttt{ArrivalDelayGroups}, \texttt{DivAirportLandings} \\
\bottomrule
\end{tabular}
\end{center}

\textbf{Enforcement.} The feature list is asserted against the banned list in code before
any model is fitted, and the train/validation/test split happens before any target
dependent statistic is computed. No test data is used to fit scalers, encoders or models.

\textbf{Transformations and their justification.}
\begin{itemize}
  \item Identifier and near duplicate columns removed: they carry no independent signal
        or duplicate \texttt{CRSElapsedTime}.
  \item Cyclic encoding of scheduled departure time: required so that midnight wraps
        correctly. Necessary for distance and margin based models (KNN, SVM) and linear
        models; harmless for trees.
  \item One hot encoding of \texttt{Airline}, \texttt{Origin}, \texttt{Dest} with a
        minimum frequency threshold, so rare airports do not each receive a column.
  \item Standardisation of numeric columns for scale sensitive models (Logistic and
        Linear Regression, KNN, SVM). Tree based models receive raw values, since they
        split on thresholds and are scale invariant.
\end{itemize}

\textbf{Sampling and splitting.} SVM and SVR scale roughly quadratically to cubically
with row count and become impractical beyond a few tens of thousands of rows. Rather
than evaluate different models on different samples, one stratified sample of
\metric{30{,}000 rows} is shared by every model so comparisons remain fair. It is split
\metric{60/20/20} into train (18,000), validation (6,000) and test (6,000), stratified
on \texttt{ArrDel15} at every step so the 19.5\% delay rate is preserved throughout. The
test set is evaluated once, for the final reported numbers only.

% ---------------------------------------------------------------- 5
\section{Models and Results}

Four algorithms were trained per problem, chosen to span different learning styles:
linear (baseline), instance based, a boosted ensemble and a margin based model.

\subsection*{Classification (target \texttt{ArrDel15})}

\begin{center}
\small
\begin{tabular}{@{}l c c c c c@{}}
\toprule
\textbf{Model} & \textbf{Train AUC} & \textbf{Val AUC} & \textbf{Test AUC} &
\textbf{Test F1} & \textbf{Gap} \\
\midrule
Logistic Regression (baseline) & 0.968 & 0.957 & \textbf{0.9649} & 0.821 & \textbf{0.011} \\
K-Nearest Neighbors            & 0.941 & 0.898 & 0.8969 & 0.559 & 0.043 \\
Gradient Boosting (XGBoost)    & 0.993 & 0.957 & 0.9631 & \textbf{0.835} & 0.036 \\
Support Vector Machine (RBF)   & 0.985 & 0.956 & 0.9639 & 0.831 & 0.030 \\
\bottomrule
\end{tabular}
\end{center}

\subsection*{Regression (target \texttt{ArrDelay}, minutes)}

\begin{center}
\small
\begin{tabular}{@{}l c c c c c@{}}
\toprule
\textbf{Model} & \textbf{Train MAE} & \textbf{Val MAE} & \textbf{Test MAE} &
\textbf{Test RMSE} & \textbf{Test $R^2$} \\
\midrule
Linear Regression (baseline)  & 7.17 & 7.35 & \textbf{7.29} & \textbf{10.35} & \textbf{0.958} \\
KNN Regressor                 & 11.83 & 12.27 & 12.54 & 18.36 & 0.868 \\
Gradient Boosting Regressor   & 6.60 & 8.40 & 8.29 & 19.18 & 0.856 \\
Support Vector Regression     & 8.14 & 9.20 & 9.39 & 31.31 & 0.615 \\
\bottomrule
\end{tabular}
\end{center}

\textbf{Why ROC-AUC and not Accuracy.} With a 80/20 split, a model predicting ``on time''
for every flight scores about 80\% Accuracy while catching zero delayed flights. ROC-AUC
is threshold independent and undistorted by the imbalance, so it is the primary
classification metric. RMSE is reported alongside MAE because it penalises large errors
more heavily, and the largest errors are the operationally important ones.

% ---------------------------------------------------------------- 6
\section{Overfitting and Hyperparameter Behaviour}

Using the rule that a train minus validation gap below 0.02 generalises well and above
0.05 is substantial, no model is severely overfit at its chosen settings. Logistic
Regression is the tightest (0.011); XGBoost is the loosest of the non-KNN models (0.036,
train AUC 0.993 against validation 0.957), indicating it has begun to fit training
noise. KNN shows the weakest performance overall, hurt by the high dimensional one hot
space where distance loses meaning.

Three hyperparameter sweeps were run, and all three produced the same shape: validation
performance rises to a peak and then declines while training performance keeps climbing.

\begin{center}
\small
\begin{tabular}{@{}l l l l@{}}
\toprule
\textbf{Model} & \textbf{Parameter} & \textbf{Behaviour observed} & \textbf{Chosen} \\
\midrule
KNN & \texttt{n\_neighbors} &
$k{=}1$ memorises: train 1.000, val 0.707, gap 0.293. Gap falls to 0.016 by $k{=}100$. &
$k=50$ \\
XGBoost & \texttt{learning\_rate} &
Validation peaks at 0.03 (0.961); gap grows 0.012 $\rightarrow$ 0.052 as rate rises to 0.40. &
0.03 \\
SVM & \texttt{C} &
Validation peaks at $C{=}1$ (0.956); gap reaches 0.060 by $C{=}100$. &
$C=1$ \\
\bottomrule
\end{tabular}
\end{center}

The lesson is consistent: model complexity must be tuned against a validation set, never
selected by what looks strongest on training data.

% ---------------------------------------------------------------- 7
\section{Model Selection}

\textbf{Selected classifier: Logistic Regression.} It attains the best test ROC-AUC
(0.9649, marginally ahead of XGBoost at 0.9631), has the smallest train--validation gap
of all four models, trains in a fraction of a second, and produces directly interpretable
coefficients, which matters for an operational tool where staff may need to know why a
flight was flagged. XGBoost achieves a slightly higher F1 (0.835 against 0.821), but
that gain does not justify the added complexity and the larger overfitting gap.

\textbf{Selected regressor: Linear Regression.} It records the best test $R^2$ and MAE
of all four models, beating every more complex alternative outright. This is coherent
with the problem structure: \texttt{ArrDelay} is close to a linear function of
\texttt{DepDelay} plus a smaller adjustment, so a linear model already captures most of
the signal, and additional flexibility costs generalisation and speed without buying
accuracy.

Selection considered performance, generalisation, interpretability and computational
cost together, not score alone. Notably, neither advanced model clearly beat its
baseline: for regression the baseline is the outright winner.

% ---------------------------------------------------------------- 8
\section{Error Analysis: Three Findings}

Confusion matrix for the selected classifier on the 6,000 row test set:
TN 4{,}502, FP 329, FN 125, TP 1{,}044.

\begin{enumerate}[leftmargin=1.5em,itemsep=3pt]
  \item \textbf{False negatives are flights that departed on time.} The 125 missed
        delays have a mean \texttt{DepDelay} of just 0.15 minutes, effectively identical
        to correctly predicted on time flights, against 71.6 minutes for true positives.
        These flights lost their time \emph{in the air}, after the moment the model makes
        its prediction. No feature available at departure could have caught them. This is
        a hard limit of the framing, not a fixable modelling weakness.
  \item \textbf{False positives are genuinely borderline.} The 329 false alarms have a
        mean \texttt{DepDelay} of 12.3 minutes and \texttt{TaxiOut} of 22.6, both clearly
        above the correct negative averages ($-3.2$ and 15.4). These are flights that
        looked risky at departure and then recovered time en route, an uncertain case
        rather than a model error.
  \item \textbf{Regression residuals are mildly but systematically biased.} For on time
        flights the mean residual is $-2.6$ minutes (slight over prediction), while for
        \emph{every} bucket of actually delayed flights it is consistently positive at
        $+4$ to $+6$ minutes. The model under predicts how severe a delay becomes once
        one is underway. The largest individual errors follow the same in air disruption
        pattern as the classification false negatives.
\end{enumerate}

Which error matters more depends on use. False negatives (a passenger told the flight is
fine when it is not) are more damaging operationally than false positives, which argues
for keeping the balanced class weighting that favours Recall on the delayed class.

% ---------------------------------------------------------------- 9
\section{Unsupervised Learning}

\textbf{Clustering.} K-Means and DBSCAN were applied to six flight shape features
(\texttt{CRSElapsedTime}, \texttt{Distance}, \texttt{dep\_sin}, \texttt{dep\_cos},
\texttt{DepDelay}, \texttt{TaxiOut}), standardised because both methods rely on
distance. \textbf{The target was never used.}

K-Means used $k{=}4$: silhouette technically peaks at $k{=}2$ (0.323) but produces one
dominant split, while $k{=}4$ (0.286) yields interpretable, differently sized segments.

\begin{itemize}
  \item \textbf{Cluster 0} (4,599): long haul, 277 min and 1,935 mi, on time on average,
        18\% delayed.
  \item \textbf{Cluster 1} (10,222): short flights departing earlier in the day, 24\%
        delayed, the worst of the three normal clusters.
  \item \textbf{Cluster 2} (14,567): short and medium flights departing later, the
        calmest group at 14\% delayed.
  \item \textbf{Cluster 3} (612): mean \texttt{ArrDelay} \metric{245.8 minutes}, 100\%
        delayed. K-Means isolated a severe delay segment entirely on its own, from
        schedule and ground features alone, without ever seeing the target.
\end{itemize}

DBSCAN used \texttt{min\_samples}~$=12$ (twice the six clustering features, a standard
rule of thumb) and \texttt{eps}~$=0.90$, chosen from the knee of the k-distance plot
(distance to each point's 12th nearest neighbour, sorted). At these settings DBSCAN does
not partition the data into several groups at all: it finds one dense core cluster
(29,299 flights, 97.7\%, mean \texttt{ArrDelay} 1.4 minutes, 18\% delayed) and marks
\metric{701 flights (2.3\%)} as noise, points too sparse to belong to any dense
neighbourhood. That noise set has mean \texttt{ArrDelay} 185.4 minutes, mean
\texttt{TaxiOut} 39.1 minutes against 16.6 in the core, and a 92\% delay rate, essentially
the same severely delayed segment K-Means isolated, reached without ever specifying how
many clusters to look for. Agreement is partial rather than exact: 52.6\% of K-Means'
612 flight cluster is labelled noise by DBSCAN, and 45.9\% of DBSCAN's noise falls inside
that K-Means cluster. The two algorithms disagree at the margin because they define
``unusual'' differently, K-Means by distance to a centroid, DBSCAN by local density.

\textbf{PCA.} One hot encoding produces \metric{220 features}. Retaining 90\% of the
variance requires \metric{44 principal components}, a five fold reduction. Logistic
Regression trained on the reduced representation scores test ROC-AUC 0.9640 against
0.9649 on the full feature set, a difference of 0.0009. PCA is therefore useful here as
\emph{compression} rather than as an accuracy improvement: valuable for a slower model
such as KNN or SVM, or a storage constrained deployment, but unnecessary for performance,
since Logistic Regression already handles 220 features without overfitting. Projected
onto the first two components the classes overlap heavily, which is expected, as those
two components capture only a modest share of total variance.

% ---------------------------------------------------------------- 10
\section{Conclusions and Recommendations}

\textbf{Most important finding.} The 0.96 correlation between departure and arrival
delay both enables this project and bounds it.

\textbf{Highest impact decision.} Fixing the leakage boundary at ``moment of departure''
rather than ``moment of landing''. That choice, made in Part 1 before any modelling,
influenced the final numbers more than any algorithm or hyperparameter selected
afterwards. A strictly pre departure model on the same data would top out far lower.

\textbf{Limitations.}
\begin{itemize}
  \item The model is only usable \emph{after} a flight has departed. It cannot address
        the harder and more valuable pre departure problem.
  \item It cannot observe weather or air traffic control conditions en route, which is
        precisely where its residual errors originate.
  \item Results are based on a 30,000 row sample of a single year (2018), constrained by
        SVM scalability.
\end{itemize}

\textbf{With more time or data.} Add live weather along the route to attack the in air
blind spot; build a strictly pre departure model side by side so the value of waiting
until departure becomes directly measurable rather than argued; and refit the selected
linear models on the full 5.58M rows, which is computationally trivial once SVM is
dropped.

\textbf{Recommendation for real world use.} Yes, for post departure updates: passenger
notification, gate assignment and connection planning, given strong and well generalising
test performance. No, for any decision that must be made before departure, which is a
different and substantially harder problem this project deliberately did not attempt.

% ---------------------------------------------------------------- 11
\section{Reflection}

The most challenging part was deciding precisely where the leakage boundary sits. It is
easy to state that no post landing column may be used, but \texttt{TaxiOut} and
\texttt{WheelsOff} sit exactly on the line and required an explicit, defensible decision
about when ``now'' is.

The most surprising result was that the baselines won. Linear Regression beat every more
complex regressor outright, and Logistic Regression matched or beat every more complex
classifier on ROC-AUC while generalising better than all of them. Additional complexity
did not translate into a better model.

On model selection, the best model is not the most sophisticated one; it is the one that
generalises reliably and fits the deployment constraints of speed and interpretability.
On overfitting, every sweep in Part 11 showed the same rise then fall in validation
performance, which reframed overfitting from a binary failure state into a matter of
degree that must be checked at every complexity setting.

If repeated, we would build the pre departure model in parallel from the start, so the
two framings could be compared directly rather than one being argued relative to the
other.

\end{document}
